\documentclass[11pt,letterpaper]{article}
\usepackage{../assets/scientific_report}
\usepackage{float}
\usepackage{listings}
\usepackage{xcolor}
\usepackage{graphicx}
\usepackage{geometry}
\geometry{headheight=23pt, top=2cm, bottom=2cm, left=2.5cm, right=2.5cm}
\setlength{\parskip}{0.5em}
\setlength{\parindent}{0pt}

\lstset{
    language=C,
    basicstyle=\ttfamily\small,
    keywordstyle=\color{blue}\bfseries,
    commentstyle=\color{green!60!black}\itshape,
    stringstyle=\color{red!80!black},
    numbers=left,
    numberstyle=\tiny\color{gray},
    numbersep=8pt,
    frame=single,
    breaklines=true,
    captionpos=b,
    tabsize=4,
    showstringspaces=false
}

\begin{document}
\pagestyle{plain}

\begin{center}
    {\LARGE\bfseries Contagem de Primos e Race Conditions}\\[0.4em]
    {\large Paralelismo de Laços e Cláusula Reduction com OpenMP}\\[1em]
    {\normalsize Eugênio Vitor Lopes dos Santos}\\[0.3em]
    {\small DCA3703 -- Programação Paralela, Universidade Federal do Rio Grande do Norte}\\[0.3em]
    {\small \today}
\end{center}


\section{Enunciado}
Implementar um programa em C que conte quantos números primos existem entre 2 e um valor máximo $N$. Depois, paralelizar o laço principal usando a diretiva \texttt{\#pragma omp parallel for} sem alterar a lógica original. Comparar o tempo de execução e os resultados das versões sequencial e paralela. Observar possíveis diferenças no resultado e no desempenho, refletindo sobre desafios como correção e distribuição de carga.

\section{Fundamentação teórica}

\subsection{Paralelismo de Laços (Loop Parallelism)}
Na programação paralela, é comum distribuir iterações independentes de um laço \texttt{for} entre diferentes threads. A diretiva \texttt{\#pragma omp parallel for} divide as iterações automaticamente. Sob o escalonamento padrão (\texttt{static}), o compilador divide os índices em blocos iguais e os distribui para as threads em tempo de compilação.

\subsection{Condição de Corrida (Race Condition)}
Uma \textit{Race Condition} surge quando múltiplas threads operam sem sincronização sobre uma mesma variável em memória compartilhada, e pelo menos uma altera o dado. 
A CPU não incrementa \texttt{count++} em um único ciclo. O hardware quebra a instrução em três passos: lê a memória para o registrador (\texttt{LOAD}), soma no registrador (\texttt{ADD}) e escreve de volta na memória (\texttt{STORE}). Se duas threads executam o \texttt{LOAD} juntas, ambas leem o mesmo número. Elas somam 1 localmente e escrevem o mesmo resultado. O programa perde um incremento e retorna um valor incorreto.

\subsection{Reduction}
Se o código travar a variável global com \texttt{critical} em cada iteração, as threads formarão uma fila e a execução ficará mais lenta que a versão sequencial. 
A diretiva \texttt{reduction(op:var)} resolve isso criando uma cópia privada da variável para cada thread no escopo do laço, com valor zero. As threads incrementam suas cópias locais na cache sem travar o barramento. Ao final do bloco paralelo, o OpenMP soma os subtotais na variável global.

\section{Experimento}

\subsection{Configuração}
O código utiliza $N=5.000.000$ elementos. A função verifica os divisores ímpares até $\sqrt{n}$. Testar números maiores consome mais ciclos de CPU. Isso desequilibra a carga de trabalho no laço.

O script executa três versões com as flags \texttt{-O3} e \texttt{-fopenmp}:
\begin{enumerate}
    \item \textbf{Sequencial}: 1 thread. Sem OpenMP.
    \item \textbf{OpenMP Naive}: Usa \texttt{\#pragma omp parallel for}. As threads modificam a variável global de contagem \texttt{count} causando corrupção de memória.
    \item \textbf{OpenMP Correct}: Usa \texttt{reduction(+:count)} protegendo o somatório.
\end{enumerate}

\section{Código-fonte}

\begin{lstlisting}[caption={tarefa05.c}]
#define _POSIX_C_SOURCE 199309L

#include <stdio.h>
#include <stdlib.h>
#include <time.h>
#include <omp.h>

#define DEFAULT_N 5000000L

static double agora(void) {
    struct timespec t;
    clock_gettime(CLOCK_MONOTONIC, &t);
    return t.tv_sec + t.tv_nsec / 1e9;
}

int is_prime(long n) {
    if (n <= 1) return 0;
    if (n == 2) return 1;
    if (n % 2 == 0) return 0;
    for (long i = 3; i * i <= n; i += 2) {
        if (n % i == 0) return 0;
    }
    return 1;
}

int main(int argc, char **argv) {
    long n = argc > 1 ? atol(argv[1]) : DEFAULT_N;
    if (n <= 0) return 1;

    int n_threads = omp_get_max_threads();
    double inicio_seq = agora();
    long count_seq = 0;
    
    for (long i = 2; i <= n; ++i) {
        if (is_prime(i)) {
            count_seq++;
        }
    }
    double tempo_seq = agora() - inicio_seq;

    double inicio_omp_naive = agora();
    long count_omp_naive = 0;

    #pragma omp parallel for
    for (long i = 2; i <= n; ++i) {
        if (is_prime(i)) {
            count_omp_naive++;
        }
    }
    double tempo_omp_naive = agora() - inicio_omp_naive;

    double inicio_omp_red = agora();
    long count_omp_red = 0;

    #pragma omp parallel for reduction(+:count_omp_red)
    for (long i = 2; i <= n; ++i) {
        if (is_prime(i)) {
            count_omp_red++;
        }
    }
    double tempo_omp_red = agora() - inicio_omp_red;

    printf("N=%ld threads=%d seq_time=%.9f omp_naive_time=%.9f omp_red_time=%.9f seq_count=%ld omp_naive_count=%ld omp_red_count=%ld\n", 
           n, n_threads, tempo_seq, tempo_omp_naive, tempo_omp_red, count_seq, count_omp_naive, count_omp_red);

    return 0;
}
\end{lstlisting}

\section{Resultados}

A execução no shell (variando de 1 a 28 threads) obteve o seguinte perfil para a busca de primos até $N=5.000.000$.

\subsection{Speedup e Observação de Race Conditions}
O programa identificou $348.513$ primos nas versões Sequencial e Reduction. A versão Naive falhou: rodando em 2 threads relatou apenas $165.441$, e em 8 threads caiu para $40.786$. Os incrementos se sobreescreveram milhares de vezes na memória.

\begin{figure}[H]
\centering
\includegraphics[width=0.9\textwidth]{speedup_plot.png}
\caption{Gráfico de Speedup para a contagem de primos com Reduction. O desbalanceamento de carga (testar números maiores demora mais) impede que a curva atinja a diagonal ideal linear (linha preta tracejada).}
\label{fig:speedup}
\end{figure}

\section{Discussão}

Este problema roda estritamente dentro da CPU e da cache L1. É totalmente CPU-bound e deve escalar. Contudo, há dois problemas mecânicos:

\subsection{Erros Silenciosos}
Falhas de memória geralmente quebram a aplicação (\textit{segmentation fault}). As Race Conditions não. O software compila, executa rápido e imprime um número no terminal. Só sabemos que o resultado está errado porque comparamos com o sequencial. É o programador quem deve garantir a corretude usando reduções locais ou estados imutáveis, em vez de compartilhar contadores nuzos.

\subsection{O Desbalanceamento de Carga (Load Imbalance)}
A Figura \ref{fig:speedup} mostra que a linha verde fica abaixo da reta tracejada de speedup perfeito. A responsabilidade disso recai sobre o escalonamento estático do compilador. 

Se o OpenMP quebra um vetor de $10.000$ iterações no meio, a Thread 0 julga os menores números e resolve tudo muito rápido. A Thread 1 fica com números maiores, roda divisões longas, e passa muito mais tempo na ALU. A Thread 0 termina o lote e fica ociosa esperando, desperdiçando potencial do processador.

\section{Conclusão}

O \texttt{\#pragma omp parallel for} expõe as threads sem que precisemos reescrever a lógica interna do for. Ainda assim, precisamos tratar a memória usando \texttt{reduction} e balancear a carga. Em problemas onde as iterações têm durações diferentes, o escalonador \texttt{static} perde eficiência. Para corrigir o problema na Figura \ref{fig:speedup}, o código pode usar escalonamento dinâmico (\texttt{schedule(dynamic, 100)}). Assim, as threads que acabam mais rápido buscam mais lotes de trabalho, evitando ociosidade.

\end{document}
